\documentclass[11pt]{article}
\usepackage[margin=2.2cm]{geometry}
\usepackage{amsmath,amssymb,booktabs,graphicx,hyperref,microtype,subcaption}
\usepackage{xcolor,placeins}
\graphicspath{{../../figures/shots_2026-09-11/}}
\hypersetup{colorlinks=true,linkcolor=blue!55!black,urlcolor=blue!55!black}

\title{Finite-shot quantum equilibrium propagation for\\antiferromagnetic phase searching}
\author{}
\date{16 September 2026}

\begin{document}
\maketitle

\begin{abstract}
We study how the number of measurements used for a quantum-equilibrium-propagation
(QEP) gradient affects training toward antiferromagnetic order.  The experiments use
open Ising chains of five and nine spins, train the squared staggered magnetisation
$\langle\mathrm{AFM}^2\rangle$ toward one, and vary the fixed shot count $N_g$.
The study separates conditional gradient noise from output-measurement noise, compares
multi-seed parameter paths, and measures performance both by epoch and by total
gradient-shot budget.  All numerical conclusions in this report are based on the
production data described in Sec.~\ref{sec:reproducibility}.
\end{abstract}

\section{Model and estimator}
\label{sec:model}

For an open chain of $N$ spins, the final system Hamiltonian is
\begin{equation}
H(\Omega,\Delta)=\frac{\Omega}{2}\sum_i\sigma_i^x
-\Delta\sum_i n_i+\sum_{i=1}^{N-1} n_i n_{i+1},
\qquad n_i=\frac{1+\sigma_i^z}{2}.
\end{equation}
The interaction has no additional factor $1/2$; this is the convention of
\texttt{qep\_shot.afm.Ising\_Chain} and of the phase diagrams used below.  We train
the order parameter
\begin{equation}
\mathrm{AFM}=\frac{1}{N}\sum_i(-1)^i\sigma_i^z,
\qquad C=\mathrm{AFM}^2,
\end{equation}
toward the target $C=1$.  The field schedules and the QEP nudge use the cosine ramp
$s(x)=[1-\cos(\pi x)]/2$ on $x\in[0,1]$, so the nudge has zero slope at both ends.
The free and nudged equilibrium states are obtained by exact diagonalisation of the
final Hamiltonians.

For a trainable parameter $\theta\in\{\Omega,\Delta\}$, the finite-shot estimator is
\begin{equation}
\widehat g_\theta(\beta)=
\frac{\overline{A}_{\theta,\beta}-\overline{A}_{\theta,0}}{\beta},
\qquad A_\theta=\partial_\theta H .
\end{equation}
Conditional on the measured output $\bar y$, independent free and nudged shot sets give
\begin{equation}
\operatorname{Var}(\widehat g_\theta\mid\bar y)=
\frac{\operatorname{Var}_\beta(A_\theta)+
      \operatorname{Var}_0(A_\theta)}{N_g\beta^2}.
\label{eq:variance}
\end{equation}
Because $A_\Omega$ and $A_\Delta$ are measured in noncommuting $X$ and $Z$ bases,
respectively, each of the free and nudged states requires an independent group of
$N_g$ shots per component.  One epoch therefore costs $4N_g$ individual gradient
measurement preparations, excluding the $N_y$ output shots used to measure the
trained observable.  An earlier version of this report (through 14 September 2026)
instead quoted $2N_gN_{\rm epoch}$, treating a free/nudged pair as one shared
``paired preparation''; Sec.~\ref{sec:solver-accuracy} documents this correction,
made together with a related ground-state-solver accuracy finding on 14--15
September 2026.  The $N=5$ figures dated 2026-09-15 and the $N=9$ figures dated
2026-09-16 discussed in Sec.~\ref{sec:performance} and Sec.~\ref{sec:budget} below
use the corrected $4N_gN_{\rm epoch}$ axis; the remaining, unreplicated $N=9$
figures (the phase panels of Sec.~\ref{sec:paths} and the expanded-seed budget
check of Fig.~\ref{fig:budget-noise-check}) still display the original
$2N_gN_{\rm epoch}$ paired-preparation axis, which is noted explicitly wherever it
appears.

Every training uses $\beta=0.1$, learning rate $\eta=0.01$, $N_y=100$ output shots,
and a fixed $N_g$; no adaptive controller is used.  Three random initial points were
drawn once with seed 20260911 and fixed at
\begin{equation}
(1.0285,2.1162),\quad(0.6955,1.9557),\quad(0.8875,1.6021).
\end{equation}
These three points, with five independent paths of 200 epochs each, underlie the
phase panels of Sec.~\ref{sec:paths} for $N=9$, the original single-fixed-start
$N=9$ budget dataset of Sec.~\ref{sec:budget}, and the measured-versus-exact
comparison of Fig.~\ref{fig:measured-exact} for both sizes.  For $N=5$
(2026-09-15) and, since 2026-09-16, for $N=9$ as well,
Sec.~\ref{sec:performance}'s performance/budget curves instead draw on an
expanded ensemble of ten initial points (these three plus seven more, drawn the
same way, shared between the two sizes) $\times$ five seeds each, described where
it is used.  For $N=5$ only, Sec.~\ref{sec:paths}'s phase panels additionally draw
on an illustrative subset of that same expanded ensemble (Fig.~\ref{fig:paths}
caption); the $N=9$ phase panels were not re-selected from it and still show the
original three points above.

\section{Results}

\subsection{Gradient distributions}

For this figure the output is always the exact ground-state expectation
$y=\langle\mathrm{AFM}^2\rangle_{\rm GS}$; it is never replaced by a finite-shot
estimate.  Only the free and nudged gradient observables are resampled.  This isolates
the direct test of Eq.~\eqref{eq:variance} from output-measurement noise.  Every
distribution is shown together with the exact shot-free expectation of the QEP
gradient at the same finite $\beta$.


At the representative point $(\Omega,\Delta)=(0.5,1.5)$ and $\beta=0.1$, the exact
ground-state outputs are 0.6705 and 0.5072 for $N=5$ and
9\footnote{Validating the new \texttt{shot\_study/exact\_solver.py}
(Sec.~\ref{sec:solver-accuracy}) against these two numbers found
$0.671697$ and $0.513241$ at the same point instead -- a $\sim$0.1--1\% gap consistent
with JAX's default complex64 precision at the time these two reference values were
first computed, not the unrelated 100-iteration training-solver bias discussed in
Sec.~\ref{sec:solver-accuracy}.}, and the exact QEP-gradient
expectations are respectively $(0.6873,0.3305)$ and $(1.7718,0.8372)$.  The sampled
histograms contract visibly onto those references.  Log--log fits of standard deviation
against $N_g$ give exponents $(-0.498,-0.488)$ for $(g_\Omega,g_\Delta)$ at $N=5$ and
$(-0.494,-0.496)$ at $N=9$, in quantitative agreement with $N_g^{-1/2}$.  For example,
the $N=9$ standard deviations shrink from $(5.99,3.81)$ at $N_g=10$ to
$(0.186,0.123)$ at $N_g=10{,}000$.

At fixed $N_g=1000$, the standard deviation grows monotonically as $\beta$ is reduced,
consistent with the $1/\beta$ amplification in Eq.~\eqref{eq:variance}; there is no
intermediate minimum in variance alone.  The total RMSE relative to the exact
$\beta=0.01$ gradient is U-shaped because large $\beta$ introduces finite-difference
bias.  Its component-wise minima occur at $(0.5,0.5)$ for $N=5$ and $(0.5,0.2)$ for
$N=9$.  This distinction between variance and total estimation error explains the
otherwise ambiguous ``optimal $\beta$'' expectation.

\begin{figure}[ht]
  \centering
  \begin{subfigure}{0.48\linewidth}\includegraphics[width=\linewidth]{gradient_distribution_ng_N5.pdf}\caption{$N=5$}\end{subfigure}
  \begin{subfigure}{0.48\linewidth}\includegraphics[width=\linewidth]{gradient_distribution_ng_N9.pdf}\caption{$N=9$}\end{subfigure}
  \caption{Finite-shot gradient distributions as $N_g$ changes.  The output in every
  sample is the exact ground-state expectation $\langle\mathrm{AFM}^2\rangle_{\rm GS}$;
  the vertical reference lines are the exact shot-free expectations of the corresponding
  finite-$\beta$ QEP gradient components.}
  \label{fig:gradient-distributions}
\end{figure}

\begin{figure}[ht]
  \centering
  \begin{subfigure}{0.48\linewidth}\includegraphics[width=\linewidth]{gradient_std_scaling_N5.pdf}\caption{$N=5$}\end{subfigure}
  \begin{subfigure}{0.48\linewidth}\includegraphics[width=\linewidth]{gradient_std_scaling_N9.pdf}\caption{$N=9$}\end{subfigure}
  \caption{Conditional gradient standard deviation versus $N_g$. Dashed lines are
  log--log fits.}
  \label{fig:gradient-scaling}
\end{figure}

\begin{figure}[ht]
  \centering
  \begin{subfigure}{0.48\linewidth}\includegraphics[width=\linewidth]{gradient_beta_tradeoff_N5.pdf}\caption{$N=5$}\end{subfigure}
  \begin{subfigure}{0.48\linewidth}\includegraphics[width=\linewidth]{gradient_beta_tradeoff_N9.pdf}\caption{$N=9$}\end{subfigure}
  \caption{Noise amplification and total finite-$\beta$ tradeoff at fixed $N_g=1000$.}
  \label{fig:beta-tradeoff}
\end{figure}

Three supplementary $\beta$ scans at $N_g=10,100,10000$ were added on 2026-09-14 to
complement the $N_g=1000$ scan above, using the same seven $\beta$ values, 256
repetitions, exact ground-state output, and fixed representative point.  Both gradient
components' sampled RMSE reaches its minimum at the largest scanned value $\beta=1$
for $N=5$ at every one of these three additional shot counts and for $N=9$ at
$N_g=10$; at $N=9$, $N_g=100$ both minima instead occur at $\beta=0.5$.  At
$N_g=10000$ both components' RMSE minimum moves inward to $\beta=0.2$ for both $N=5$
and $N=9$.  A boundary minimum at $\beta=1$ does not locate a true optimum beyond the
scanned range; only the $N_g=1000$ and $N_g=10000$ scans show an interior minimum for
both sizes.

\begin{figure}[ht]
  \centering
  \begin{subfigure}{0.32\linewidth}\includegraphics[width=\linewidth]{gradient_beta_tradeoff_N5_Ng10.pdf}\caption{$N=5$, $N_g=10$}\end{subfigure}
  \begin{subfigure}{0.32\linewidth}\includegraphics[width=\linewidth]{gradient_beta_tradeoff_N5_Ng100.pdf}\caption{$N=5$, $N_g=100$}\end{subfigure}
  \begin{subfigure}{0.32\linewidth}\includegraphics[width=\linewidth]{gradient_beta_tradeoff_N5_Ng10000.pdf}\caption{$N=5$, $N_g=10000$}\end{subfigure}
  \begin{subfigure}{0.32\linewidth}\includegraphics[width=\linewidth]{gradient_beta_tradeoff_N9_Ng10.pdf}\caption{$N=9$, $N_g=10$}\end{subfigure}
  \begin{subfigure}{0.32\linewidth}\includegraphics[width=\linewidth]{gradient_beta_tradeoff_N9_Ng100.pdf}\caption{$N=9$, $N_g=100$}\end{subfigure}
  \begin{subfigure}{0.32\linewidth}\includegraphics[width=\linewidth]{gradient_beta_tradeoff_N9_Ng10000.pdf}\caption{$N=9$, $N_g=10000$}\end{subfigure}
  \caption{Supplementary $\beta$ noise--bias tradeoff scans at $N_g=10,100,10000$
  (Fig.~\ref{fig:beta-tradeoff} shows $N_g=1000$), same 256-repetition
  exact-ground-state-output protocol.}
  \label{fig:beta-tradeoff-supp}
\end{figure}

\FloatBarrier
\subsection{From random walks to smooth parameter paths}
\label{sec:paths}


The geometric transition is clear for both sizes.  For $N=5$, increasing $N_g$ from
10 to 10,000 reduces the median path-length/endpoint ratio from 11.39 to 1.10 and
raises the mean cosine between adjacent steps from $-0.006$ to 0.693.  The corresponding
$N=9$ values change from 11.47 to 1.22 and from 0.032 to 0.496.  Thus low-shot paths
are locally almost uncorrelated random walks, whereas high-shot replicas from the same
initial point nearly collapse onto a smooth common curve.

\begin{figure}[ht]
  \centering
  \begin{subfigure}{0.48\linewidth}\includegraphics[width=\linewidth]{phase_trajectories_N5_Ng10.pdf}\caption{$N=5$, $N_g=10$}\end{subfigure}
  \begin{subfigure}{0.48\linewidth}\includegraphics[width=\linewidth]{phase_trajectories_N5_Ng10000.pdf}\caption{$N=5$, $N_g=10000$}\end{subfigure}
  \begin{subfigure}{0.48\linewidth}\includegraphics[width=\linewidth]{phase_trajectories_N9_Ng10.pdf}\caption{$N=9$, $N_g=10$}\end{subfigure}
  \begin{subfigure}{0.48\linewidth}\includegraphics[width=\linewidth]{phase_trajectories_N9_Ng10000.pdf}\caption{$N=9$, $N_g=10000$}\end{subfigure}
  \caption{Five sampling-seed paths from each of three starting points, overlaid on
  the exact ground-state $\langle\mathrm{AFM}^2\rangle$ landscape; every line is one
  actual seed and no mean path is drawn.  For $N=9$ these are the original three
  randomly drawn points of Sec.~\ref{sec:model}.  For $N=5$ they are instead a
  performance-informed illustrative selection (starts labelled 3 and 5, plus a
  replacement start 1) drawn from the ten-point ensemble of
  Sec.~\ref{sec:performance} and chosen on 2026-09-14 for visual separation and for
  reaching the ordered region; they are not a statistical sample and do not replace
  the ten-start ensemble averages reported in Sec.~\ref{sec:performance}.}
  \label{fig:paths}
\end{figure}

\FloatBarrier
\subsection{Performance by epoch and physical shot count}
\label{sec:performance}

\paragraph{$N=5$ and $N=9$ (expanded ten-start ensembles, 2026-09-15 and 2026-09-16).}
The panels of Fig.~\ref{fig:performance} and the main $N=5$/$N=9$ panels of
Fig.~\ref{fig:budget} now use an identical expanded-ensemble protocol at both sizes:
the ten initial points of Sec.~\ref{sec:model} $\times$ five sampling seeds per point
(50 paths per $N_g$), replacing the original three-start/five-seed ensemble used for
performance and budget figures through 14 September 2026 ($N=5$) and 15 September
2026 ($N=9$).  Because $N_g=10,100,1000,10000$ visit very different physical-shot
ranges per epoch, they are trained for different lengths -- 5000, 500, 200, and 200
epochs respectively -- chosen so all four curves cover a common cumulative-shot
window of roughly $40{,}000$--$200{,}000$ gradient shots under the corrected
$4N_gN_{\rm epoch}$ convention.  Shaded bands in both sizes' panels are $\pm1$ SEM
across the ten start-means, each already averaging five seeds, and are not a 95\%
interval.  At each curve's own final epoch the mean measured order is, for $N=5$,
$0.887\pm0.018$ ($N_g=10$, epoch 5000), $0.950\pm0.009$ ($N_g=100$, epoch 500),
$0.802\pm0.077$ ($N_g=1000$, epoch 200), and $0.798\pm0.081$ ($N_g=10000$, epoch
200); for $N=9$ (completed 2026-09-16, an independent 200-case Raven campaign
replicating the $N=5$ protocol exactly, Sec.~\ref{sec:reproducibility}),
$0.745\pm0.040$ ($N_g=10$), $0.898\pm0.037$ ($N_g=100$), $0.718\pm0.117$
($N_g=1000$), and $0.712\pm0.120$ ($N_g=10000$).  $N=9$'s final measured order is
lower than $N=5$'s at every $N_g$ (by 0.04--0.14) and carries roughly twice the SEM;
this is consistent with $N=9$'s previously documented greater training sensitivity
and initial-condition dependence (see ``The larger chain exposes optimization
variability'' below) rather than a new anomaly or a convention mismatch, since both
ensembles now share an identical measurement and uncertainty convention.  As for
$N=5$, because the higher-$N_g$ curves are capped at only 200 epochs, $N=9$'s
ensemble also does not show a monotonic epoch-count benefit of larger $N_g$; both
sizes instead reflect the joint effect of gradient noise and the update budget,
consistent with the fixed-budget results of Sec.~\ref{sec:budget}.  Both ensembles,
at both sizes, retain the historical approximate training-state solver of
Sec.~\ref{sec:solver-accuracy}.  With this addition, the expanded-ensemble
performance/budget comparison that was previously available only for $N=5$ is now
complete for both system sizes.

\begin{figure}[ht]
  \centering
  \begin{subfigure}{0.48\linewidth}\includegraphics[width=\linewidth]{performance_vs_epoch_N5.pdf}\caption{epochs, $N=5$ (10 starts $\times$5 seeds)}\end{subfigure}
  \begin{subfigure}{0.48\linewidth}\includegraphics[width=\linewidth]{performance_vs_gradient_shots_N5.pdf}\caption{$4N_gN_{\rm epoch}$ shots, $N=5$}\end{subfigure}
  \begin{subfigure}{0.48\linewidth}\includegraphics[width=\linewidth]{performance_vs_epoch_N9.pdf}\caption{epochs, $N=9$ (10 starts $\times$5 seeds)}\end{subfigure}
  \begin{subfigure}{0.48\linewidth}\includegraphics[width=\linewidth]{performance_vs_gradient_shots_N9.pdf}\caption{$4N_gN_{\rm epoch}$ shots, $N=9$}\end{subfigure}
  \caption{Mean measured order parameter with uncertainty bands over all sampled paths
  per $N_g$.  Both sizes now use the identical expanded ten-start$\times$five-seed
  ensemble, the same corrected $4N_gN_{\rm epoch}$ physical-shot axis, and the same
  band convention: $\pm1$ SEM across the ten start-means, each already averaging five
  seeds; see text.}
  \label{fig:performance}
\end{figure}

\begin{figure}[ht]
  \centering
  \begin{subfigure}{0.48\linewidth}\includegraphics[width=\linewidth]{performance_measured_vs_exact_N5.pdf}\caption{$N=5$}\end{subfigure}
  \begin{subfigure}{0.48\linewidth}\includegraphics[width=\linewidth]{performance_measured_vs_exact_N9.pdf}\caption{$N=9$}\end{subfigure}
  \caption{Measured (dashed, $N_y=100$ output shots) versus exact free ground-state
  (solid) $\langle\mathrm{AFM}^2\rangle$ at the same pre-update parameters, averaged
  over the original three starts $\times$ five seeds for both sizes (this comparison
  was not extended to either size's expanded ten-start ensemble).  Exact values use a float64 dense diagonalization (\texttt{scipy.linalg.eigh})
  computed outside the notebook in \texttt{shot\_study/trajectory\_expectations.py},
  not the approximate training-state solver of Sec.~\ref{sec:solver-accuracy}.  Epochwise RMSE between the
  two curves is $0.0073$--$0.0078$ for $N=5$ and $0.0050$--$0.0052$ for $N=9$.  At the
  last plotted (pre-update) epoch the exact means are approximately
  $(0.7111,0.9009,0.9204,0.9205)$ for $N=5$ and $(0.6089,0.6710,0.7365,0.7085)$ for
  $N=9$, for $N_g=10,100,1000,10000$ respectively.}
  \label{fig:measured-exact}
\end{figure}

\FloatBarrier
\subsection{Fixed total-shot budgets}
\label{sec:budget}

The reference Fixed\_TN budget values $5{,}000$, $10{,}000$, $20{,}000$, $50{,}000$,
and $100{,}000$ are used for the original single-fixed-start datasets at both
$N=9$ and $N=5$, described immediately below; for each $N_g$ the number of
completed updates is $\lfloor N_{g,{\rm budget}}/(2N_g)\rfloor$ under that
(paired-preparation) convention, and points with zero possible updates are left
blank rather than assigned an artificial initial performance.  Both original
datasets are superseded, as main-panel figures, by the expanded ten-start
ensembles described further below.

\paragraph{Original single-fixed-start datasets, $N=9$ and $N=5$ (historical,
superseded).}
The optimal fixed shot count is consistently intermediate rather than maximal.  For
this $N=5$ dataset (single fixed start $(\Omega,\Delta)=(0.5,1.5)$, five seeds, the
full grid $N_g\in\{10,\dots,10000\}$), the best $N_g$ at budgets
$(5,10,20,50,100)\times10^3$ is respectively $(20,20,50,20,200)$, with mean exact
orders $(0.925,0.935,0.988,0.989,0.964)$.  For $N=9$ the corresponding optima are
$(50,50,100,200,100)$ and $(0.946,0.980,0.942,0.974,0.979)$.  None of these ten
comparisons selects either $N_g=10$ or $N_g\ge500$: too few shots make each update
unreliable, whereas too many consume the budget before enough parameter updates can
be made.  The nonmonotonic movement of the empirical optimum with budget should be
interpreted at the resolution of five training seeds, rather than as a precisely
determined scaling law.  Neither dataset is any longer the one plotted as the main
budget panel in Fig.~\ref{fig:budget} below: $N=5$'s was superseded 2026-09-15 and
$N=9$'s 2026-09-16, both by the expanded ten-start ensembles described in the next
two paragraphs.  The $N=5$ dataset is additionally retained in the single-start
25-seed robustness check further below (Fig.~\ref{fig:budget-noise-check-n5}); the
$N=9$ dataset has no analogous 25-seed full-grid check, only the narrower
$N_g\in\{20,50,100,200\}$, $S=50{,}000$ expanded-seed check of
Fig.~\ref{fig:budget-noise-check} below.

\paragraph{$N=5$ main panel (expanded ten-start ensemble, 2026-09-15, corrected
$4N_g$ convention).}
Using the same ten-start$\times$five-seed ensemble as Sec.~\ref{sec:performance}, the
left panel of Fig.~\ref{fig:budget} instead shows exact ground-state
$\langle\mathrm{AFM}^2\rangle$ at physical gradient-shot caps $S=10{,}000$,
$20{,}000$, $40{,}000$, $100{,}000$, and $200{,}000$ (corrected $4N_g$ convention;
these correspond to the former paired caps $B=S/2$), for only
$N_g=10,100,1000,10000$ -- the shot counts actually simulated for this ensemble.
Completed updates are $k=\lfloor S/(4N_g)\rfloor$; cases with $k=0$ are omitted.  Here
the ranking is essentially reversed from the single-start dataset above: the best
$N_g$ is $10$ at $S=10{,}000,20{,}000,40{,}000$ (means $0.741,0.779,0.865$) and $100$
at $S=100{,}000,200{,}000$ (means $0.859,0.947$, versus $0.839,0.886$ for $N_g=10$);
$N_g=1000$ and $10000$ remain at or below $0.51$ at every budget.  This reversal is
consistent with a different initial-condition distribution -- the ten-start
ensemble's mean initial $\langle\mathrm{AFM}^2\rangle$ is 0.370 (range
0.172--0.863) versus 0.672 for the single fixed start -- combined with the
per-$N_g$ epoch caps of Sec.~\ref{sec:performance}: at these budgets, $N_g=1000$ and
$10000$ have completed at most 50 and 5 updates respectively, too few from a harder
average starting point, whereas $N_g=10,100$ have completed enough updates to
approach their own late-epoch performance.  This is not a claim that $N_g\ge1000$ is
generally worse; it reflects this particular budget/epoch-cap combination and
initial-condition sample, not a re-derivation of the single-start optimum above.

\paragraph{$N=9$ main panel (expanded ten-start ensemble, 2026-09-16, corrected
$4N_g$ convention).}
Using the identical shared ten-point$\times$five-seed ensemble (the same ten
initial points as $N=5$, Sec.~\ref{sec:model}) and the same protocol as the $N=5$
main panel just described, the right panel of Fig.~\ref{fig:budget} shows the
analogous $N=9$ result at the same physical gradient-shot caps
$S=10{,}000,20{,}000,40{,}000,100{,}000,200{,}000$ and the same
$N_g=10,100,1000,10000$.  The ranking is qualitatively similar to $N=5$ but the
crossover to the more accurate estimator is delayed: $N_g=10$ is best at
$S=10{,}000,20{,}000,40{,}000,100{,}000$ (means $0.676,0.682,0.744,0.762$) and only
$N_g=100$ overtakes it at $S=200{,}000$ (mean $0.898$ versus $0.769$ for
$N_g=10$); $N_g=1000$ and $10000$ again remain at or below $0.42$ at every budget.
For $N=5$ the crossover to $N_g=100$ instead already occurs at $S=100{,}000$.  This
delayed crossover is consistent with $N=9$'s ten-start ensemble having an even
lower and harder mean initial $\langle\mathrm{AFM}^2\rangle=0.241$ (range
$0.090$--$0.829$, versus $N=5$'s $0.370$, range $0.172$--$0.863$), computed by
exact diagonalization at the same ten shared points: from a harder average start,
$N_g=100$ needs correspondingly more of the fixed shot budget before its
lower-noise updates catch up with $N_g=10$'s larger number of noisier ones.  As
for $N=5$, this is not a claim that $N_g\ge1000$ is generally worse, only a
description of this particular budget/epoch-cap combination and initial-condition
sample.  With this panel, the fixed-budget comparison of the expanded ensemble --
previously available only for $N=5$ -- is now complete for both system sizes.

\begin{figure}[ht]
  \centering
  \begin{subfigure}{0.48\linewidth}\includegraphics[width=\linewidth]{time_budget_N5.pdf}\caption{$N=5$: 10 starts $\times$5 seeds, $S=4N_gN_{\rm epoch}$ (2026-09-15)}\end{subfigure}
  \begin{subfigure}{0.48\linewidth}\includegraphics[width=\linewidth]{time_budget_N9.pdf}\caption{$N=9$: 10 starts $\times$5 seeds, $S=4N_gN_{\rm epoch}$ (2026-09-16)}\end{subfigure}
  \caption{Exact order at fixed total gradient-shot budgets, for the identical
  expanded ten-start$\times$five-seed ensemble and the identical corrected
  $4N_gN_{\rm epoch}$ convention at both sizes.  $N=9$'s curves sit systematically
  below $N=5$'s at matched $(N_g,S)$, and its crossover to $N_g=100$ occurs at a
  larger budget; see text for the initial-condition explanation.  The
  original single-fixed-start datasets and legacy $2N_gN_{\rm epoch}$ convention
  discussed above are no longer plotted here (for $N=9$'s original dataset, see the
  Fig.~\ref{fig:budget-noise-check} check below instead; for $N=5$'s, see
  Fig.~\ref{fig:budget-noise-check-n5}).}
  \label{fig:budget}
\end{figure}
\FloatBarrier

\paragraph{$N=9$ expanded-seed check (original single-fixed-start dataset,
legacy convention).}
The conspicuous fall and recovery of the $N=9$, $50{,}000$-shot curve at $N_g=50$
--from the original single-fixed-start dataset described two paragraphs above, not
the expanded ten-start ensemble of Fig.~\ref{fig:budget}-- was checked with 20
additional independent seeds at each of $N_g\in\{20,50,100,200\}$.  The original five-seed means over these four points
were $(0.953,0.869,0.948,0.974)$, whereas the combined 25-seed means are
$(0.900,0.927,0.962,0.952)$.  Thus the deep, isolated $N_g=50$ dip is not
reproduced: the expanded curve rises smoothly through $N_g=100$ and then decreases
slightly.  Its 95\% confidence intervals are respectively
$[0.861,0.940]$, $[0.898,0.957]$, $[0.946,0.979]$, and $[0.935,0.970]$.
Welch tests give $p=0.261$ for $N_g=50$ versus 20 and $p=0.143$ versus 200.
The $N_g=100$ mean is modestly higher than the $N_g=50$ mean
($\Delta=0.0348$, unadjusted two-sided $p=0.0389$), so the verification supports
a broad intermediate optimum but not the original sharp V-shaped feature.

\begin{figure}[ht]
  \centering
  \includegraphics[width=0.62\linewidth]{time_budget_N9_noise_check.pdf}
  \caption{Expanded-seed check at $N=9$ and a total gradient-shot budget of
  $50{,}000$.  Blue points show combined 25-seed means with 95\% confidence
  intervals; the dashed grey line is the original five-seed mean.}
  \label{fig:budget-noise-check}
\end{figure}

\FloatBarrier
\paragraph{$N=5$ single-start 25-seed check (2026-09-14).}
The apparent dips of the original five-seed, single-fixed-start
$(\Omega,\Delta)=(0.5,1.5)$ budget curves (legacy paired budget $B=50{,}000$ dip at
$N_g=50$; $B=100{,}000$ dip at $N_g=100$) were checked the same way, with 20
additional seeds at every $N_g$ and budget.  The $B=50{,}000$, $N_g=50$ mean moves
from $0.910216$ (5 seeds) to $0.958448$ (new 20 seeds) to $0.948801$ (combined 25);
its log-$N_g$-neighbour interpolation contrast shrinks from $-0.068188$ to
$-0.018076$ (95\% percentile-bootstrap interval $[-0.046503,0.008196]$).  The
$B=100{,}000$, $N_g=100$ mean moves from $0.915755$ to $0.973736$ to $0.962140$; its
contrast changes from $-0.033573$ to $+0.015520$ (95\% interval
$[-0.004986,0.034153]$), i.e.\ that dip disappears entirely.  Both results support
sampling fluctuation in the original five-seed curves, not a corrected optimum and
not a claim of monotonicity: a small $B=50{,}000$ dip remains at the 25-seed point
estimate (Fig.~\ref{fig:budget-noise-check-n5}), though its neighbour-contrast
interval still includes zero.  This single-start dataset and its historical solver
are distinct from -- and were never mixed into -- the ten-start ensemble of the main
$N=5$ panel above.

\begin{figure}[ht]
  \centering
  \includegraphics[width=0.62\linewidth]{time_budget_N5_25seeds.pdf}
  \caption{Single-fixed-start $(\Omega,\Delta)=(0.5,1.5)$, 25-seed fixed-budget check
  for $N=5$ (legacy paired-budget convention, full $N_g\in\{10,\dots,10000\}$ grid).
  Bands are pointwise bootstrap 95\% intervals.  This is the dataset described in the
  preceding paragraph, distinct from the ten-start ensemble of Fig.~\ref{fig:budget}.}
  \label{fig:budget-noise-check-n5}
\end{figure}

\FloatBarrier
\subsection{Why adaptive shot counts become attractive}


Over the last 50 epochs of the $N_g=1000$ diagnostic, the conditional gradient means
are $(0.0260,0.00450)$ for $N=5$ and $(0.0133,0.00166)$ for $N=9$, while the
corresponding finite-shot standard deviations remain $(0.497,0.0991)$ and
$(0.669,0.105)$.  In particular, the absolute $g_\Omega$ noise is roughly five to six
times the $g_\Delta$ noise and does not vanish as the signal approaches zero.  Median
late relative-noise ratios are $(21.1,24.3)$ and $(54.1,70.7)$, explaining the violent
fluctuations of a fixed-shot estimator near convergence and motivating component-wise
adaptive shot allocation.

\begin{figure}[ht]
  \centering
  \begin{subfigure}{0.48\linewidth}\includegraphics[width=\linewidth]{gradient_moments_during_training_N5.pdf}\caption{$N=5$}\end{subfigure}
  \begin{subfigure}{0.48\linewidth}\includegraphics[width=\linewidth]{gradient_moments_during_training_N9.pdf}\caption{$N=9$}\end{subfigure}
  \caption{Conditional gradient signal, finite-shot standard deviation, relative
  noise, and measured training output for a typical $N_g=1000$ run.}
  \label{fig:moments}
\end{figure}

\FloatBarrier
\section{Discussion and observations}

\paragraph{Gradient noise is the dominant controlled effect.}
Figure~\ref{fig:gradient-distributions} deliberately removes output-shot noise by using
the exact ground-state value of $\langle\mathrm{AFM}^2\rangle$ in every call.  The
resulting agreement with the exact gradient and the fitted $N_g^{-1/2}$ scaling show
that the broad low-shot distributions are intrinsic to sampling the gradient
observables, rather than an artefact of a fluctuating target signal.  The training runs
restore the experimentally relevant $N_y=100$ output measurement, so their scatter
contains both sources of noise.

\paragraph{Smoothness and efficiency are different questions.}
Large $N_g$ is unambiguously beneficial when cost is measured in epochs: paths become
smoother and replicas become more reproducible.  It is not unambiguously beneficial
when cost is measured in experiments, and exactly which $N_g$ wins depends on the
ensemble.  For the original single-fixed-start datasets (both $N=9$ and $N=5$), the
fixed-budget optima sit at $N_g=20$--200: an extremely accurate update can be less
useful than several moderately noisy updates, and the optimum moves nonmonotonically
across the five budgets.  For the expanded ten-start ensembles at both sizes
(Sec.~\ref{sec:performance}--\ref{sec:budget}), by contrast, $N_g=10$ or $100$
dominate at every tested budget and $N_g\ge1000$ never wins, because those runs are
additionally capped at only 200 epochs from a harder average starting point.
Within the expanded ensembles, $N=9$'s crossover from $N_g=10$ to $N_g=100$ is
delayed to a larger shot budget than $N=5$'s ($S=200{,}000$ versus $S=100{,}000$),
consistent with $N=9$'s even lower and harder mean initial
$\langle\mathrm{AFM}^2\rangle$ (0.241 versus 0.370, Sec.~\ref{sec:budget}) rather
than any difference in methodology.  The qualitative lesson -- an intermediate
$N_g$ beats the extremes under a fixed shot budget -- is robust across all four
ensembles, but the specific optimal value is not a universal constant; it depends
on the initial-condition distribution and the epoch budget actually allotted to
each $N_g$, not only on the physics of the estimator.

\paragraph{The larger chain exposes optimization variability.}
The $N=9$ paths become geometrically smoother as $N_g$ grows, but the between-start
spread of their final performance remains large.  In the expanded ten-start
ensemble, $N=9$'s final-epoch mean measured order is lower than $N=5$'s at every
$N_g$ (e.g.\ $0.745\pm0.040$ versus $0.887\pm0.018$ at $N_g=10$) and carries roughly
twice the SEM (Sec.~\ref{sec:performance}), and $N_g=1000$ again has a slightly
higher late mean than $N_g=10000$ ($0.718\pm0.117$ versus $0.712\pm0.120$), although
the difference is far below either quantity's own SEM across the ten starts.  Thus
suppressing estimator noise cannot by itself eliminate initial-condition sensitivity
in the more structured landscape, and this sensitivity now shows up consistently in
both the original and the expanded $N=9$ ensembles.

\paragraph{Shot allocation should be component and time dependent.}
Near convergence the mean gradients approach zero, while the absolute $g_\Omega$ noise
stays roughly five to six times larger than the $g_\Delta$ noise.  A single fixed shot
count therefore over-resolves some updates and under-resolves others.  This is the most
direct numerical motivation in the study for allocating shots separately by parameter
and increasing them only when the signal-to-noise ratio requires it.

The numerical scope should nevertheless be kept in view: the distribution study uses
one parameter point, the original trajectory averages cover three starts and five
seeds rather than an asymptotic ensemble (ten starts for the expanded $N=5$ and,
since 2026-09-16, $N=9$ figures), and, as the next subsection details, state
preparation is an approximate iterative solver, not an exact diagonalization as its
name suggests.

\subsection{Ground-state solver accuracy}
\label{sec:solver-accuracy}

Every gradient and every trajectory in this report -- $N=5$ and $N=9$, original and
expanded -- prepares the free and nudged states with \texttt{QEP.grad.Exact\_Diagonalization},
which despite its name delegates to \texttt{QEP.model.get\_gs\_sparse}: a fixed
100-iteration momentum-gradient power iteration from a uniform initial state, with no
residual convergence check (\texttt{QEP/model.py}).  An independent SciPy float64
\texttt{eigh} diagonalization of the same saved Hamiltonians at several checkpoints
found several-percent-level gradient-mean discrepancies, e.g.\ at the first saved
point for $N=5$: saved (100-iteration) mean $(0.759542,0.364020)$ versus exact
$(0.752488,0.357070)$; for $N=9$: saved $(1.675777,0.793209)$ versus exact
$(1.606324,0.757700)$.  A direct convergence check comparing 100/200/500 iterations
against the same float64 reference gives maximum component-relative mean errors of
$1.9460\%/0.01696\%/0.000991\%$ at the first $N=5$ checkpoint,
$3.0509\%/0.03704\%/0.03887\%$ at the last $N=5$ checkpoint (not monotonically
improving: the late-training small-$\Delta$ signal remains sensitive to float32-level
numerical precision even at 500 iterations), and $4.6841\%/0.004943\%/0.000121\%$ at
the first $N=9$ checkpoint.  A separate check found the cross-component correlation
of independently sampled $g_\Omega,g_\Delta$ to be $-0.001614$ (95\% interval
$[-0.04493,0.04170]$), ruling out shared-randomness key reuse as an alternative
explanation for the near-identical relative-noise curves of Fig.~\ref{fig:moments};
this is a genuine state-preparation accuracy issue, not a sampling artefact (full
details in \texttt{shot\_study/figure13\_noise\_audit\_2026-09-14.md}).

This bias affects every trajectory, gradient distribution, and performance/budget
curve in this report through their state preparation.  It does \emph{not} affect the
independently recomputed exact-reference quantities that this report also uses:
phase-diagram backgrounds, the endpoint values of Fig.~\ref{fig:paths}, the solid
curves of Fig.~\ref{fig:measured-exact}, and both panels of Fig.~\ref{fig:budget}
($N=5$ since 2026-09-15 and $N=9$ since 2026-09-16, the latter via the shared,
validated \texttt{exact\_solver.py})
were all evaluated by genuine post-hoc dense diagonalization of the saved trajectory
parameters, not by this iterative solver -- but those exact endpoints are evaluated
\emph{on top of} parameter paths that were themselves trained with the approximate
solver, so they do not certify that the training itself was exact.  We have not
corrected the training/gradient-sampling solver itself: a fully corrected retraining
campaign is deferred pending a scope decision on how much compute to spend re-running
the study and how to avoid mixing corrected-solver trajectories with the existing
approximate-solver dataset used throughout this report.  A validated, documented,
reusable exact solver (\texttt{shot\_study/exact\_solver.py}, wrapping the same
float64 \texttt{scipy.linalg.eigh} lowest-eigenpair approach already used ad hoc in
several analysis scripts) is now available for that future work and for exact
endpoint evaluation; it is deliberately not wired into the training loop and has not
been used to alter any existing trajectory or production data.

\section{Conclusions}

The experiments establish that finite-shot QEP gradients concentrate on their exact
finite-$\beta$ expectations at the predicted $N_g^{-1/2}$ rate.  Increasing $N_g$
turns random-walk-like training into smooth motion and improves performance per epoch,
but an intermediate $N_g$ is superior under fixed measurement budgets in every
ensemble tested -- $N_g=20$--200 for $N=9$ and $N=5$'s original single-start
datasets, $N_g=10$--100 for the expanded ten-start ensembles now available at both
sizes -- rather than either extreme.  Within the expanded ensembles, $N=9$
consistently trails $N=5$ in final measured order and delays its own crossover to
$N_g=100$ to a larger shot budget, consistent with $N=9$'s previously documented
greater optimization sensitivity and its harder mean initial condition, not a
methodological difference between the two sizes (Secs.~\ref{sec:performance}
and~\ref{sec:budget}).  Because the residual noise is strongly component dependent
and eventually dominates the decaying gradient signal, these fixed-shot results
support adaptive allocation by both parameter and training stage.  Separately, every
result in this report inherits a quantified, several-percent-level state-preparation
bias from the approximate 100-iteration ground-state solver used throughout
(Sec.~\ref{sec:solver-accuracy}); a corrected rerun is future work, not a silent
correction applied here.  With the 2026-09-16 $N=9$ campaign, the expanded-ensemble
performance and fixed-budget comparison between $N=5$ and $N=9$ is now complete at
both system sizes, closing out a comparison previously available only for $N=5$.

\section{Reproducibility and data integrity}
\label{sec:reproducibility}

The self-contained implementation is in \texttt{AFM\_searching/afm\_shots}; experiment
drivers, the central configuration, validation, plotting code, and Slurm files are in
\texttt{AFM\_searching/shot\_study}.  Raw production data are compressed NPZ files in
\texttt{data/shots\_2026-09-11}; every file contains JSON metadata with the physical
settings, seed, and whether it is a smoke or production run. Reduced local tests carry
an explicit \texttt{\_smoke} suffix and are excluded from analysis.  The Raven run used
  jobs 30179398, 30177302, 30177303, and 30177304 in the dedicated directory
\texttt{/ptmp/qinwang/notebooks/QEP\_fixed\_shot\_20260911}.  The validator checks all
226 expected production files, metadata consistency, finite arrays, and the intentional
zero-update mask before plotting.  The focused budget-noise verification used Zeropoint
array 8895092 and replacement element 8895181, producing 80 additional production files
(20 seeds for each of $N_g=20,50,100,200$).  All 80 retained runs completed with exit
code 0; one non-progressing allocation was intentionally cancelled and rerun on another
node.  The additional files and their Slurm logs are stored alongside the original data
and logs, and the reduced statistics are recorded in
\texttt{budget\_noise\_check\_summary.json}.

The expanded $N=5$ ensemble (Secs.~\ref{sec:performance}--\ref{sec:budget}) was trained
in a separately isolated Raven root,
\texttt{/raven/u/qinwang/notebooks/QEP\_overlap\_20260914/QEP\_multiple\_shot}, and
summarized locally by \texttt{shot\_study/summarize\_overlap.py} into
\texttt{trajectory\_overlap\_summary\_N5.npz}, which validates completion, metadata,
shapes, and exact preservation of the original three-start prefixes before any
figure is drawn from it.  The single-fixed-start 25-seed budget check
(Fig.~\ref{fig:budget-noise-check-n5}) was produced on Zeropoint arrays 8895474 and
8895517 and combined locally by \texttt{shot\_study/figure11b\_resampling.py}.  A
validated, reusable exact ground-state solver for post-hoc analysis --
\texttt{shot\_study/exact\_solver.py}, discussed in Sec.~\ref{sec:solver-accuracy} --
is checked against the exact references quoted above for $N=5,9$ at
$(\Omega,\Delta)=(0.5,1.5)$; it is not used by the production training/sampling
pipeline described above, and none of the original 226 production files (or the
subsequent expanded/verification files) were regenerated or altered while producing
this report revision.

The expanded $N=9$ ensemble (Secs.~\ref{sec:performance}--\ref{sec:budget}) was
added on 2026-09-16, replicating the $N=5$ campaign's protocol exactly (the same
10-point$\times$5-seed$\times N_g\in\{10,100,1000,10000\}$ design, the same
5000/500/200/200 epoch schedule, and the same reused original three-point
prefixes) via 200 new cases submitted as Raven job 30255973 (25 groups of 8, capped
at 8 concurrent GPUs) in the same isolated Raven root used for $N=5$.  All 25
groups completed with exit code \texttt{0:0} and empty error logs; a comprehensive
log grep for tracebacks, CUDA errors, out-of-memory events, and non-finite values
found nothing.  The 200 resulting production files were summarized locally (the
same \texttt{summarize\_overlap.py}, run locally rather than on the cluster because
the remote conda environment lacked \texttt{threadpoolctl}, mirroring the $N=5$
precedent) into \texttt{trajectory\_overlap\_summary\_N9.npz}, with a maximum
eigenpair residual of $5.76\times10^{-14}$, well below the $10^{-9}$ validation
gate.  \texttt{summarize\_overlap.py} was refactored during this pass to call the
shared \texttt{exact\_solver.py} instead of its own duplicated diagonalization
code; this refactor was regression-tested against the existing, already-delivered
$N=5$ summary (bit-identical arrays, max absolute difference $0.0$) before being
relied on for the new $N=9$ summary.  Figures 9, 10, and 11b of
\texttt{figs\_N9.ipynb} were updated to use this expanded ensemble, the corrected
$4N_gN_{\rm epoch}$ convention, and exact-endpoint diagonalization, matching
\texttt{figs\_N5.ipynb}; no other cell in either notebook was changed.  Full
details are recorded in \texttt{GPT\_relay.md} under ``2026-09-16: N=9 expanded
ensemble (Figures 9, 10, 11b) replicating the N=5 campaign'' and the closing
dated entry that follows it.

\end{document}
